\clearpage\fancyhead[L]{\small PALLAS.}\hypertarget{native-361-lower}{}\label{native-361-lower}
\begin{center}\textbf{PALLAS.}\par\vspace{4pt}\hrule\vspace{4pt}\fontsize{10}{12}\selectfont Göttingische gelehrte Anzeigen. Stück 127. Seite 1257\;..\;1259. 1812 Aug. 8.\par\vspace{4pt}\hrule\vspace{4pt}\end{center}

\par\indent Im 67. Stück dieser Blätter [S. 349 d. B.] ist bereits die erste Auffindung der Pallas in diesem
Jahre und die schöne Uebereinstimmung ihrer Stellung mit der nach der Störungstheorie berechneten
Ephemeride angezeigt. Die Anzahl der seitdem auf der hiesigen Sternwarte angestellten Beobachtungen
ist nur klein, und besonders verstattete um die Zeit der Opposition ungünstiges Wetter, Mondschein und
der hohe, für das Local unbequeme, Stand des Planeten nur eine einzige. Die sämmtlichen in diesem
Jahre vom Prof. \textsc{Gauss} auf unserer Sternwarte gemachten Beobachtungen sind folgende:

\begin{center}
\fontsize{10}{12}\selectfont\setlength{\tabcolsep}{3pt}\renewcommand{\arraystretch}{1.15}\begin{tabular}{c|c|c|c}
1812 & Mittlere Zeit & \shortstack{Gerade\\Aufsteigung} & \shortstack{Nordliche\\Abweichung} \\
\hline
April 9 & $11^{\mathrm u}\;32^{\mathrm m}\;27^{\mathrm s}$ & $269^\circ\;5'\;58''4$ & $16^\circ\;38'\;0''$ \\
Mai 2 & 10 \quad 36 \quad 6 & 268 \quad 42 \quad 5.3 & 21 \quad 4 \quad 31.7 \\
— 4 & 10 \quad 44 \quad 44 & 268 \quad 31 \quad 38.0 & 21 \quad 25 \quad 4.3 \\
Juni 7 & 11 \quad 13 \quad 51 & 262 \quad 50 \quad 51.1 & 25 \quad 2 \quad 56.6 \\
\end{tabular}
\end{center}

\par\indent Bei der letzten Beobachtung hatte der Planet die zehnte Grösse.

\par\indent Prof. \textsc{Gauss} hatte mehrere seiner auswärtigen astronomischen Freunde, denen bessere Hülfsmittel
zu Gebote stehen, zur Beobachtung der Pallas um die Zeit der Opposition eingeladen: den bereits eingegangenen
Nachrichten zufolge ist dies indess weder auf der kaiserl. Sternwarte in Paris, noch in
Marseille und Berlin geglückt: man darf sich darüber nicht wundern, da der dies Jahr so lichtschwache
Planet an den fixen Instrumenten kaum die geringste Beleuchtung verträgt. Blos Hr. \textsc{Burckhardt} war
so glücklich, auf der Sternwarte der Militärschule in Paris folgende Meridianbeobachtungen zu erhalten:

\begin{center}
\fontsize{10}{12}\selectfont\setlength{\tabcolsep}{3pt}\renewcommand{\arraystretch}{1.15}\begin{tabular}{c|c|c|c}
1812 & Mittlere Zeit & \shortstack{Gerade\\Aufsteigung} & \shortstack{Nordliche\\Abweichung} \\
\hline
Juni 5 & $12^{\mathrm u}\;35^{\mathrm m}\;57^{\mathrm s}4$ & $263^\circ\;15'\;24''0$ & $24^\circ\;58'\;28''5$ \\
6 & 12 \quad 33 \quad 10.5 & 263 \quad 2 \quad 35 & 25 \quad 0 \quad 30 \\
8 & 12 \quad 21 \quad 37.0 & 262 \quad 37 \quad 6.5 & 25 \quad 3 \quad 57.7 \\
11 & 12 \quad 7 \quad 15.7 & 261 \quad 58 \quad 37.8 & 25 \quad 6 \quad 29.3 \\
\end{tabular}
\end{center}

\par\indent Prof. \textsc{Gauss} übergab diese Beobachtungen Hrn. \textsc{Nicolai}, um daraus die Opposition zu berechnen.
Sein Resultat ist folgendes:
\SourceLeaf{362}\begin{center}\fontsize{10}{12}\selectfont\setlength{\tabcolsep}{0pt}\renewcommand{\arraystretch}{1.15}\begin{tabular}{l}
Opposition der Pallas 1812 Juni 10. $3^{\mathrm u}\;2^{\mathrm m}\;48^{\mathrm s}$\\mittl. Zeit im Meridian von Göttingen
\end{tabular}\end{center}
\begin{center}\fontsize{10}{12}\selectfont\setlength{\tabcolsep}{2pt}\renewcommand{\arraystretch}{1.15}\begin{tabular}{lcr}
Wahre Länge & $.\;.\;.\;.\;.\;.\;.\;.\;.\;.\;.\;.\;.\;.\;.\;.\;.\;.$ & $259^\circ\;37'\;51''7$\\
Wahre geocentrische Breite & $.\;.\;.\;.\;.\;.\;.\;.\;.\;.\;.\;.\;.$ & $48\;16\;26.1\;\mathrm{N}.$\\
\end{tabular}\end{center}

\par\indent Dies ist nunmehr die \emph{achte} beobachtete Opposition der Pallas seit ihrer
Entdeckung. Von der vor einem Jahre von Hrn. \textsc{Nicolai} im voraus geführten
Rechnung weicht das Resultat um 29 Minuten in der Zeit, 1 Minute in der
Länge, und 10 Secunden in der Breite ab (s. Gött. gel. Anz. 1811 S. 1292
[S. 329 d. B.]). Prof. \textsc{Gauss} glaubt in diesen, obwohl geringen,
Unterschieden schon ganz sicher den Einfluss der bisher noch nicht
untersuchten Störungen durch den Saturn und Mars zu erkennen: zur
Berechnung der erstern sind bereits alle Einleitungen getroffen, und wir
werden davon zu seiner Zeit umständlicher Rechenschaft geben. Einstweilen
hat Prof. \textsc{Gauss} untersucht, wie genau die sämmtlichen acht bisher
beobachteten Oppositionen sich noch darstellen lassen, wenn blos die
Störungen durch den Jupiter in Betracht gezogen werden. Das Resultat ist,
dass man, um die möglichste Uebereinstimmung zu erhalten,

\begin{center}\fontsize{9}{11}\selectfont\setlength{\tabcolsep}{2pt}\renewcommand{\arraystretch}{1.15}\begin{tabular}{lcr}
\shortstack[l]{die Epoche der mittlern Länge\\am 10. Juni 1812 um} & $.\;.\;.\;.$ & $-26''95$\\
die tägliche mittlere Bewegung um & $.\;.\;.\;.\;.\;.\;.\;.\;.\;.$ & $-0''01225$\\
die Logarithm. der halben Axe um & $.\;.\;.\;.\;.\;.\;.\;.\;.\;.\;.$ & $+0.0000046$\\
die Länge des Perihels um & $.\;.\;.\;.\;.\;.\;.\;.\;.\;.\;.\;.\;.\;.$ & $-18''84$\\
den Excentricitätswinkel um & $.\;.\;.\;.\;.\;.\;.\;.\;.\;.\;.\;.\;.$ & $-3''92$\\
die Länge des Knoten um & $.\;.\;.\;.\;.\;.\;.\;.\;.\;.\;.\;.\;.\;.$ & $+0''80$\\
die Neigung der Bahn um & $.\;.\;.\;.\;.\;.\;.\;.\;.\;.\;.\;.\;.\;.$ & $-1''21$\\
\end{tabular}\end{center}

\Needspace{16\baselineskip}\noindent ändern müsse. Die Uebereinstimmung der beobachteten und berechneten
Opposition ist dann folgende:

\begin{center}
\fontsize{10}{12}\selectfont\setlength{\tabcolsep}{7pt}\renewcommand{\arraystretch}{1.15}\begin{tabular}{c|c|c}
 & \multicolumn{2}{c}{Unterschied im} \\
Opposition & mittlere Länge & heliocentr. Breite \\
\hline
1803 & $+$\ 6$''$4  & $-$\ 2$''$8 \\
1804 & $-$10.9 & $-$9.7 \\
1805 & $-$13.8 & $-$0.7 \\
1807 & $+$10.2 & $-$5.8 \\
1808 & $+$19.1 & $+$11.7 \\
1809 & $+$16.1 & $-$1.5 \\
1811 & $+$14.2 & $-$9.5 \\
1812 & $-$22.3 & $+$6.1 \\
\end{tabular}
\end{center}

\par\indent Die \emph{neunte} Opposition wird, nach Hrn. \textsc{Nicolai}'s Rechnung, eintreten
1813 Aug. 18. $9^{\mathrm u}$, in $325^\circ\;24'$ Länge, $24^\circ\;37'$
N. geocentr. Breite.
